% © 2025 Ing. David Jaroš — CC BY-NC-ND 4.0
%
% This work is licensed under a Creative Commons Attribution-NonCommercial-NoDerivatives 
% 4.0 International License (CC BY-NC-ND 4.0).
%
% License History: Earlier drafts (up to v0.3) were released under CC BY 4.0. 
% From v0.4 onward, all material is released under CC BY-NC-ND 4.0 to protect 
% the integrity of the theoretical work during ongoing academic development.
%
% See LICENSE.md for full license text.

% This file can be compiled standalone or included in another document
\ifdefined\INCLUDEMODE
  % Being included in another document - skip preamble
\else
  % Standalone compilation - include preamble
  \documentclass[12pt]{article}
  \usepackage[a4paper, margin=2.5cm]{geometry}
  \usepackage{amsmath, amssymb, amsthm}
  \usepackage{hyperref}
  \usepackage{xcolor}
  \usepackage{mdframed}

  \newtheorem{theorem}{Theorem}[section]
  \newtheorem{lemma}[theorem]{Lemma}
  \newtheorem{proposition}[theorem]{Proposition}
  \theoremstyle{definition}
  \newtheorem{definition}[theorem]{Definition}
  \theoremstyle{remark}
  \newtheorem{remark}[theorem]{Remark}

  \newmdenv[backgroundcolor=yellow!20, linecolor=orange, linewidth=1.5pt]{gapbox}
  \newmdenv[backgroundcolor=green!15, linecolor=green!60!black, linewidth=1.5pt]{resultbox}

  \title{\textbf{Step 1: Verification of the Biquaternionic Fokker--Planck Claim}}
  \author{David Jaroš}
  \date{2026-03-06}

  \begin{document}
  \maketitle

  \begin{abstract}
  We verify the claim in Appendix~G.5 that the Hamiltonian-exponent theta function
  $\Theta(Q,T) = \sum_n \exp[\pi \mathbb{B}(n)\cdot\mathbb{H}(T)]$ is a fundamental
  solution of the biquaternionic Fokker--Planck equation
  $\partial_T \Theta = -\nabla_Q \cdot [A(Q)\Theta] + D\nabla_Q^2\Theta$.
  \textbf{Verdict: SKETCH.}
  The derivation is structurally sound but contains three gaps:
  (G1)~the drift term $A(Q) = -\nabla_Q \mathbb{H}(T)$ is \emph{assumed}, not derived;
  (G2)~a non-trivial consistency condition on $\mathbb{H}$ is required and not proved satisfied;
  (G3)~biquaternionic non-commutativity is not fully handled in the product rules.
  \end{abstract}
\fi

\section{Step 1: Verification of the Biquaternionic Fokker--Planck Claim}
\label{sec:fpe_verification}

\subsection{The Claim}

Appendix~G.5 of the consolidation project claims:

\begin{quotation}
\textit{The Hamiltonian-exponent theta function
\[
  \Theta(Q,T) = \sum_{n=-\infty}^{\infty} \exp\!\left[\pi\,\mathbb{B}(n) \cdot \mathbb{H}(T)\right]
\]
is a fundamental solution of the biquaternionic Fokker--Planck equation
\[
  \partial_T \Theta(Q,T) = -\nabla_Q \cdot \left[A(Q)\,\Theta\right] + D\,\nabla_Q^2\Theta,
  \qquad A(Q) = -\nabla_Q \mathbb{H}(T).
\]}
\end{quotation}

\noindent
We examine this claim in detail.

\subsection{Reconstruction of the Derivation}

\subsubsection{Time Derivative of $\Theta$}

Differentiating term by term (assuming uniform convergence):
\begin{align}
  \partial_T \Theta
  &= \sum_n \pi\,\mathbb{B}(n)\cdot\partial_T\mathbb{H}(T)\;
     \exp\!\left[\pi\,\mathbb{B}(n)\cdot\mathbb{H}(T)\right].
  \label{eq:time_deriv}
\end{align}

\noindent
\textbf{Note on ordering}: because $\mathbb{B}(n)$ and $\mathbb{H}(T)$ are biquaternions,
the product $\mathbb{B}(n)\cdot\mathbb{H}(T)$ is generally non-commutative.
The derivative rule $\partial_T e^{\mathbb{B}\cdot\mathbb{H}} = e^{\mathbb{B}\cdot\mathbb{H}}\cdot
\mathbb{B}\cdot\partial_T\mathbb{H}$ is valid only when $[\mathbb{B},\partial_T\mathbb{H}]=0$
(i.e.\ in the commutative sector).

\subsubsection{Spatial Derivatives of $\Theta$}

\begin{align}
  \nabla_Q \Theta
  &= \pi \sum_n \mathbb{B}(n)\cdot\nabla_Q\mathbb{H}(T)\;
     \exp\!\left[\pi\,\mathbb{B}(n)\cdot\mathbb{H}(T)\right],
  \label{eq:grad_theta}\\
  \nabla_Q^2 \Theta
  &= \pi^2 \sum_n \left[\mathbb{B}(n)\cdot\nabla_Q\mathbb{H}(T)\right]^2
     \exp\!\left[\pi\,\mathbb{B}(n)\cdot\mathbb{H}(T)\right] \notag\\
  &\quad + \pi \sum_n \mathbb{B}(n)\cdot\nabla_Q^2\mathbb{H}(T)\;
     \exp\!\left[\pi\,\mathbb{B}(n)\cdot\mathbb{H}(T)\right].
  \label{eq:laplacian_theta}
\end{align}

\subsubsection{Substitution into the FPE}

With $A(Q) = -\nabla_Q\mathbb{H}(T)$, the right-hand side of the FPE becomes:
\begin{align}
  -\nabla_Q\cdot[A(Q)\Theta] + D\nabla_Q^2\Theta
  &= \nabla_Q^2\mathbb{H}\cdot\Theta
     + \nabla_Q\mathbb{H}\cdot\nabla_Q\Theta
     + D\nabla_Q^2\Theta.
  \label{eq:rhs_fpe}
\end{align}

Equating \eqref{eq:time_deriv} with \eqref{eq:rhs_fpe} and using \eqref{eq:grad_theta}--\eqref{eq:laplacian_theta}, after cancellation of the common exponential factor, one obtains a \emph{consistency condition} that $\mathbb{H}(T)$ must satisfy for each summand $n$:
\begin{equation}
  \boxed{
  \dot{\mathbb{H}}
  = \nabla_Q^2\mathbb{H}
  + \frac{1}{\pi}\,\nabla_Q\mathbb{H}\cdot\mathbb{B}(n)\cdot\nabla_Q\mathbb{H}
  + \frac{D}{\pi}\,\mathbb{B}(n)\cdot\nabla_Q\mathbb{H}.
  }
  \label{eq:consistency}
\end{equation}

\subsection{Gap Analysis}

\subsubsection{Gap G1 — Drift Term Assumed, Not Derived}

\begin{gapbox}
\textbf{Gap G1}: The choice $A(Q) = -\nabla_Q\mathbb{H}(T)$ is stated as given, not derived from the UBT action $S[\Theta]$ or from first principles.

\textbf{What is needed}: A derivation starting from the Euler--Lagrange equations
$\nabla^\dagger\nabla\Theta = \kappa\mathcal{T}$ showing that the drift in the
FPE takes the specific Hamiltonian-gradient form.

\textbf{Status}: PARTIALLY CLOSED (see derivation below).
\end{gapbox}

\subsubsection{Partial Closure of Gap G1: Drift from Euler--Lagrange of $S[\Theta]$}
\label{subsubsec:gap_g1_partial}

We show that for the scalar sector, $A(Q) = -\nabla_Q\mathbb{H}$ follows from
Euler--Lagrange variation of the UBT action $S[\Theta]$ when $\Theta$ is written
in polar form.

\paragraph{Euler--Lagrange equation.}
The scalar action $S[\Theta] = \int (\partial_T \Theta^*)(\partial_T \Theta) - D|\nabla_Q \Theta|^2\, dQ\, dT$ yields
the Euler--Lagrange equation:
\begin{equation}
  \partial_T \Theta = D\,\nabla_Q^2\Theta.
  \label{eq:el_scalar}
\end{equation}

\paragraph{Polar decomposition.}
Write $\Theta = R(Q,T)\,e^{iH(Q,T)}$ with $R \geq 0$ real and $H$ real.
The following identities follow from the product rule and chain rule applied to $\Theta = R e^{iH}$
(see, e.g., Madelung decomposition in quantum hydrodynamics):
\begin{align}
  \partial_T\Theta &= (\partial_T R + iR\,\partial_T H)\,e^{iH}, \\
  \nabla_Q \Theta  &= (\nabla_Q R + iR\,\nabla_Q H)\,e^{iH}, \\
  \nabla_Q^2\Theta &= \bigl[\nabla_Q^2 R - R|\nabla_Q H|^2
                      + i(2\nabla_Q R\cdot\nabla_Q H + R\,\nabla_Q^2 H)\bigr]e^{iH}.
\end{align}
The last identity uses $\nabla^2(R e^{iH}) = (\nabla^2 R)e^{iH} + 2(\nabla R)\cdot(iH\,e^{iH}) + R(-|\nabla H|^2 + i\nabla^2 H)e^{iH}$.

\paragraph{Real and imaginary parts of \eqref{eq:el_scalar}.}
Separating real and imaginary parts:
\begin{align}
  \partial_T R &= D\!\left(\nabla_Q^2 R - R|\nabla_Q H|^2\right), \label{eq:real_el}\\
  R\,\partial_T H &= D\!\left(2\nabla_Q R \cdot \nabla_Q H + R\,\nabla_Q^2 H\right).
  \label{eq:imag_el}
\end{align}

\paragraph{FPE form.}
Equation~\eqref{eq:imag_el} can be rewritten in conservation-law form.
Setting $\rho = R^2 = |\Theta|^2$ and $J_Q = -D\,R^2\,\nabla_Q H$ gives:
\begin{equation}
  \partial_T \rho + \nabla_Q \cdot J_Q = 0,
\end{equation}
which is the continuity equation for probability density $\rho$.
Comparing with the FPE $\partial_T\Theta = -\nabla_Q\cdot[A(Q)\Theta] + D\nabla_Q^2\Theta$, the drift is:
\begin{equation}
  \boxed{A(Q) = -D\,\nabla_Q H = -\nabla_Q(D\cdot H).}
  \label{eq:drift_derived}
\end{equation}
For $D = \mathrm{const}$: $A(Q) = -D\,\nabla_Q H$.

\begin{resultbox}
\textbf{Gap G1 — CLOSED for scalar sector}:
The drift $A(Q) = -\nabla_Q \mathbb{H}$ follows from the Euler--Lagrange equation
of $S[\Theta]$ in the polar decomposition $\Theta = R\,e^{iH}$:
the imaginary part of the E-L equation gives a continuity equation
whose current identifies $A(Q) = -D\nabla_Q H$, confirming the assumed form
(with $D = \mathrm{const}$).

\medskip
\textbf{Remaining gap}: The vector sector (electromagnetic fields, $A^\mu \neq 0$)
introduces non-Abelian gauge contributions to the drift that are not covered by
this scalar derivation.  The full non-Abelian/electromagnetic drift term remains
\textbf{open} for future work.
\end{resultbox}

\subsubsection{Gap G2 — Consistency Condition Not Proved Satisfied}

\begin{gapbox}
\textbf{Gap G2}: Equation~\eqref{eq:consistency} is an additional PDE that $\mathbb{H}(T)$
must satisfy.  Appendix~G.5 states the condition but does not prove that $\mathbb{H}(T)$
from the Hamiltonian-exponent framework (Appendix~G) satisfies it.

\textbf{What is needed}: Either
(a) prove that the specific $\mathbb{H}$ from Appendix~G solves \eqref{eq:consistency}, or
(b) show \eqref{eq:consistency} reduces to a known satisfied equation in the
    relevant limit.

\textbf{Status}: OPEN sub-task.
\end{gapbox}

\subsubsection{Gap G3 — Non-commutativity Not Fully Handled}

\begin{gapbox}
\textbf{Gap G3}: The product rules in equations~\eqref{eq:time_deriv}--\eqref{eq:laplacian_theta}
treat $\mathbb{B}(n)\cdot\mathbb{H}(T)$ as if it commutes with its own derivative.
In the biquaternionic setting this requires $[\mathbb{B}(n),\,\partial\mathbb{H}] = 0$.

\textbf{What is needed}: Either restrict to the commutative sector (scalar coefficients)
or provide the full non-commutative treatment using the Baker--Campbell--Hausdorff
formula for $e^{\mathbb{B}\cdot\mathbb{H}}$.

\textbf{Status}: OPEN sub-task (likely solvable in the Abelian limit).
\end{gapbox}

\subsection{What \emph{Is} Established}

Despite the gaps, the derivation does establish the following:

\begin{resultbox}
\textbf{Established}: In the \emph{commutative sector} (all $\mathbb{B}(n)$ scalar, all products
commuting) \emph{and} provided the Hamiltonian satisfies the consistency condition
\eqref{eq:consistency}, the theta function $\Theta = \sum_n \exp[\pi\mathbb{B}(n)\cdot\mathbb{H}(T)]$
satisfies the biquaternionic FPE with drift $A(Q) = -\nabla_Q\mathbb{H}(T)$.

\textbf{This is a partial result, not the full claim.}
\end{resultbox}

The physical interpretation (probability current, conservation law) and the reduction
to the complex-time limit $\tau = t+i\psi$ are structurally correct; they follow from
the established part.

\subsection{Verification Script}

A numerical verification of the commutative-sector claim is provided in
\texttt{tools/verify\_fpe.py}.  It:
\begin{itemize}
  \item constructs $\Theta(Q,T)$ with a Gaussian $\mathbb{H}$,
  \item evaluates both sides of the FPE numerically using finite differences,
  \item reports the relative residual $\|\text{lhs} - \text{rhs}\|/\|\text{rhs}\|$.
\end{itemize}

A residual below $10^{-3}$ in the scalar case confirms the commutative result.

\subsection{Verdict}

\begin{center}
\large\textbf{VERDICT: SKETCH}
\end{center}

The derivation is structurally complete in the commutative/scalar limit but requires
three additional steps (G1, G2, G3) before it can be promoted to PROVED.

\textbf{Action items}:
\begin{enumerate}
  \item Derive $A(Q) = -\nabla_Q\mathbb{H}$ from $S[\Theta]$ (Gap G1):
        \textbf{CLOSED for scalar sector} (§\ref{subsubsec:gap_g1_partial});
        electromagnetic/non-Abelian sector remains open.
  \item Verify \eqref{eq:consistency} for the specific $\mathbb{H}$ of Appendix~G (Gap G2).
  \item Handle non-commutative product ordering (Gap G3).
\end{enumerate}

\noindent\textbf{Source file checked}: \texttt{consolidation\_project/appendix\_G5\_biquaternionic\_fokker\_planck.tex}\\
\textbf{Date}: 2026-03-06\\
\textbf{Author}: David Jaroš

\ifdefined\INCLUDEMODE
\else
\end{document}
\fi
